\section{Generalized Collinear Parton Observables}
\label{sec:gpo}
%------------------------------------------------------------------------------
In the previous section, we have extensively discussed the leading-twist collinear PDFs
that characterize the 1D structure of the proton in longitudinal momentum space.
There exist various other parton observables that provide complementary information.
In this section, we focus on observables defined by collinear parton correlators, in the sense
that only the collinear quark and gluon mode contribute, corresponding to the so-called
{\it collinear expansion} in QCD factorizations~\cite{Sterman:1994ce,Collins:2011zzd}.
We call them ``generalized collinear parton observables'' (GCPOs), and discuss their
calculations through LaMET framework. For observables defined by parton correlators
involving transverse separations, in particular, the TMDPDFs, Wigner functions,
and LFWFs, we will consider them in the following sections.

One of the important GCPOs is the GPDs introduced in~\cite{Mueller:1998fv}, and
rediscovered~\cite{Ji:1996ek} from their connection to the spin structure of the proton.
%
They describe the correlation between the transverse position and longitudinal momentum of partons inside the proton, and thus provide important information for 3D imaging of the proton. A proton spin sum rule was derived in terms of the moments of the GPDs, which has
stimulated considerable general interest in the GPDs. It was also found that
in the so-called zero skewness limit or when the longitudinal momentum transfer
vanishes, the GPD has a probability interpretation in the impact parameter space~\cite{Burkardt:2000za}.
%
In general case, it is related to the quantum phase-space distributions or Wigner functions~\cite{Ji:2003ak,Belitsky:2003nz}.
Experimentally, the GPDs can be measured through hard exclusive processes such as deeply virtual Compton scattering (DVCS) or meson production (DVMP) that were first proposed in~\cite{Ji:1996ek,Ji:1996nm}. Much effort has been devoted to measuring such processes at completed and ongoing experiments, including HERA, COMPASS and JLab. For a more comprehensive discussion on the GPDs, we refer the readers to the review articles~\cite{Ji:1998pc,Ji:2004gf,Diehl:2003ny,Belitsky:2005qn}.
%
Despite that the GPDs have more complicated kinematic dependence and relation to experimental observables,
various fitting methods have been proposed in the literature to fit available DVCS and DVMP data~\cite{Kumericki:2016ehc,Favart:2015umi}. In parallel, one can also extract certain information on the GPDs from lattice calculations of their moments~\cite{Hagler:2007xi,Gockeler:2003jfa,Alexandrou:2019ali}, which, however, is again very limited due to the same difficulties existing in lattice calculations of the PDF moments. For JLab 12 GeV program and future EIC, it
is critically important to have first-principle calculations of GPDs with much better understanding
of the physical landscape in different kinematic variables.


A simpler but closely related GCPO is the parton distribution amplitudes (DAs), which are collinear matrix elements of light-cone operators between a hadron state and the QCD vacuum, representing the probability amplitude of finding a given Fock state in the hadron. They can be probed in certain exclusive processes~\cite{Brodsky:2002st}, and are crucial inputs for processes relevant to measuring fundamental parameters of the Standard Model and probing new physics. There exists a vast amount of literature on this subject, particularly about the pion DA. For a review see e.g. \cite{Brodsky:1989pv,Grozin:2005iz,Braun:2006hn}.

Another type of GCPO is the higher-twist parton distributions.
They are defined by multi-parton correlation functions, and quantify the proton structure
in terms of longitudinal momentum correlations~\cite{Jaffe:1982pm,Ellis:1982cd,Jaffe:1991ra}.
Although physically interesting, they are hard to
separate theoretically due to mixing with the
leading-twist ones~\cite{Mueller:1984vh,Ji:1994md}, and difficult to
extract experimentally because they are power-suppressed~\cite{Ji:1993ey}.
Higher-twist effects can become important in kinematic regions where the
suppression is relaxed. Moreover, some twist-three distributions,
$g_T$ and $h_L$, are different; they have no leading-twist to mix with and are dominant
in spin-related observables~\cite{Jaffe:1991ra}.
Twist-three GPDs are also relevant for studying parton OAM in the
proton~\cite{Hatta:2012cs,Ji:2012ba,Courtoy:2013oaa} and can be
accessed through DVCS process~\cite{Penttinen:2000dg,Kiptily:2002nx}.


In principle, all the GCPOs discussed above can be computed within LaMET.
%
In addition, an accurate LaMET expansion for the leading-twist PDFs requires
calculations of quasi higher-twist matrix elements.
In the following, we begin with the flavor non-singlet quark GPDs and hadronic DAs for which the computational procedure has been well established, and then give some generic discussions on higher-twist distributions, followed by the discussion on power-suppressed contributions required to extract the leading-twist quark PDFs, which have been investigated using different approaches though not yet implemented in numerical computations.


%------------------------------------------------------------------------------
\subsection{Generalized Parton Distributions}
\label{sec:others-gpd}
%------------------------------------------------------------------------------

The operators defining the GPDs are the same as those defining the PDFs. Thus, the LaMET calculation of PDFs can be rather straightforwardly generalized to the GPDs by taking into account the non-forward kinematics~\cite{Liu:2018tox}. To illustrate how it works, let us take the nonsinglet unpolarized quark GPDs in the nucleon as an example.


The unpolarized quark GPDs are defined through the following matrix element~\cite{Ji:2004gf}
\begin{align}\label{eq:GPD}
	F&=\frac{1}{2{\bar P}^+}\int \frac{d\lambda}{2\pi}e^{-i x\lambda}\langle P'S'|O_{\gamma^+}(\lambda n)|PS\rangle\nn\\
	&=\frac{1}{2{\bar P}^+}\bar{u}(P'S')\bigg[H\gamma^+
	+E\frac{i\sigma^{+\mu}\Delta_\mu}{2M}\bigg]u(PS)\,,
\end{align}
where we have suppressed the arguments $(x,\xi,t, \mu)$ of $F$, $H$ and $E$ for simplicity. The operator
\begin{align}
	O_{\gamma^+}(\lambda n)=\bar\psi(\frac{\lambda n} 2)\gamma^+ W(\frac{\lambda n}{2},-\frac{\lambda n}{2})\psi(-\frac{\lambda n} 2)
\end{align}
with $n^\mu=1/\sqrt 2(1/{\bar P}^+, 0,0, -1/{\bar P}^+)$ is the same operator used to define the unpolarized quark PDF, $M$ is the nucleon mass.
%
The momentum fraction $x\in [-1,1]$, and
\begin{align}\label{eq:kinematic_variable}
	\Delta\equiv P'-P,\;\; t\equiv \Delta^2,\;\; \xi\equiv -\frac{P'^+-P^+}{P'^++P^+}=-\frac{\Delta^+}{2{\bar P}^+}\,,
\end{align}
where without loss of generality we have chosen a Lorentz frame in which the average momentum takes the following form
\begin{align}
	{\bar P}^\mu\equiv\frac{P'^\mu+P^\mu}{2}=({\bar P}^0,0,0,{\bar P}^z)\,.
\end{align}
The skewness parameter $\xi\in[-1,1]$ since $P^+, P'^+\ge0$. Besides, there exists another kinematic constraint on $\xi$, which follows from $\vec \Delta^2_\perp\ge 0$,
\beq\label{eq:sklimit}
\xi\le\xi_{\rm max}(t)=\sqrt{\frac{-t}{-t+4M^2}}\,.
\eeq
In the following, we will also assume $\xi>0$ without loss of generality. With these kinematic constraints, the GPDs can be divided into several kinematic regions that have different physical interpretations. As shown in \fig{gpdkin}, in the region $\xi<x<1\, (-1<x<-\xi)$ the distribution describes the emission and reabsorption of a quark (antiquark), while in the region $-\xi<x<\xi$ it represents the creation of a quark and antiquark pair. The first region is similar to that present in usual PDFs and referred to as the DGLAP region, whereas the second is similar to that in a meson DA, which will be discussed later in this section, and referred to as the Efremov-Radyushkin-Brodsky-Lepage (ERBL) region. The easiest way to see this is in light-cone quantization and light-cone gauge where the matrix element defining the GPDs can be rewritten in terms of parton creation and annihilation operators, for details see e.g.~\cite{Ji:2004gf}.

\begin{figure}
	\centering
	\includegraphics[width=.9\linewidth]{images/gpd-gpd_fig09.pdf}
	\caption{Parton interpretation of the GPDs in different kinematic regions.}
	\label{fig:gpdkin}
\end{figure}

The quark GPDs defined above have a number of remarkable properties, see, e.g.,~\cite{Ji:1998pc,Ji:2004gf,Diehl:2003ny,Belitsky:2005qn}
, which either hold or have similar counterparts for the quark quasi-GPDs to be defined below. Apart from their physical significance, these properties also serve as useful checks on calculations related to GPDs.


According to LaMET, the unpolarized quark GPDs defined above can be determined by calculating the following quasi-GPDs
\begin{align}\label{eq:quasiGPD}
	\tilde F&=\frac{1}{2{\bar P}^0}\int \frac{d\lambda}{2\pi}e^{i y\lambda }\langle P'S'|{O}_{\gamma^0}(z)|PS\rangle\nn\\
	&=\frac{1}{2{\bar P}^0}\bar{u}(P'S')\bigg\{\tilde H\gamma^0
	+\tilde E\frac{i\sigma^{0\mu}\Delta_\mu}{2M}\bigg\}u(PS)\,,
\end{align}
where we have again suppressed the arguments $(y,\tilde{\xi},t,{\bar P}^z, \mu)$ of $\tilde F$, $\tilde H$, and $\tilde E$. The operator ${O}_{\gamma^0}(z)=\bar\psi(\frac{z}{2})\gamma^0 W(\frac z 2,-\frac z 2)\psi(-\frac z 2)$ is the same operator defining the unpolarized quark quasi-PDF, and $\lambda=z {\bar P}^z$.
%
As in the quasi-PDF case, the momentum fraction $y$ extends from $-\infty$ to $\infty$. The skewness parameter for the quasi-GPD
\begin{align}
	\tilde{\xi}=-\frac{P'^z-P^z}{P'^z+P^z}=-\frac{\Delta^z}{2{\bar P}^z} =  \xi +\mathcal{O}\left({M^2\over ({\bar P}^z)^2},{t\over ({\bar P}^z)^2}\right)
\end{align}
differs from the light-cone skewness $\xi$ by power suppressed corrections. Moreover, the constraint from $\vec\Delta_\perp^2\ge0$ becomes~\cite{Ji:2015qla}
\begin{align}
	{\tilde \xi}\le\frac{1}{2{\bar P}^{z}}\sqrt{\frac{-t\left[({\bar P}^{z})^{2}+M^{2}-t/4\right]}{M^{2}-t/4}},
\end{align}
which differs from the constraint in \eq{sklimit} by corrections of $\mathcal O(M^2/({\bar P}^z)^2, t/({\bar P}^z)^2)$. We can replace $\tilde \xi$ with $\xi$ and attribute the difference to generic power suppressed contributions.


The quasi-GPDs defined above can be renormalized by observing that their UV divergence depends only on the operators defining them, but not on the external states. Since ${O}_{\gamma^0}(z)$ is multiplicatively renormalized, we can choose the same renormalization factor as that for the quasi-PDF~\cite{Stewart:2017tvs,Liu:2018uuj} to renormalize the quasi-GPD.
%
After renormalization, the quasi-GPD can then be matched to the usual GPD through a factorization formula.

The factorization of quasi-GPDs was first proposed and verified at one-loop order in~\cite{Ji:2015qla,Xiong:2015nua}, where a transverse momentum cutoff and a quark mass were used as the UV and IR regulator, respectively. Later on, a detailed derivation based on OPE was given in~\cite{Liu:2019urm}. In contrast with the OPE for the quasi-PDF, a crucial difference here is that the total derivative of operators can come into play, as it simply gives momentum transfer factors when sandwiched between non-forward external states, and therefore is non-vanishing. In other words, the local twist-two operators as those in Eq.~(\ref{eq:ope-tq}) will mix under renormalization with operators with total derivatives.
%
The RGE that governs the mixing reads~\cite{Braun:2003rp},
\begin{align} \label{eq:rge}
	&\mu^2 {d\over d\mu^2}O^{{\mu}_0 {\mu}_1 \ldots {\mu}_n}(\mu) = \sum_{m=0}^{[n/2]}\Gamma_{nm}\\
	&\times\left[i\partial^{(\mu_1}\cdots i\partial^{\mu_{2m}} \bar{\psi} \gamma^{\mu_0}i\overleftrightarrow{D}^{\mu_{2m+1}}\cdots i\overleftrightarrow{D}^{\mu_n)}\psi-{\rm trace}\right]\,,\nn
\end{align}
where $\Gamma_{nm}$ is the anomalous dimension of the associated operators, $\overleftrightarrow{D}=(\overrightarrow{D}-\overleftarrow{D})/2$ with $\overrightarrow{D} (\overleftarrow{D})$ denoting the covariant derivative acting to the right (left). The above equation can be diagonalized by choosing an appropriate operator basis. Such an operator basis has been studied in the literature and known as the ``renormalization group improved'' conformal operators~\cite{Braun:2003rp,Mueller:1993hg}. In terms of the matrix elements of these operators, we have
\begin{align}\label{eq:ope-twist2}
	&\langle P'|{O}_{\gamma^0}(z)|P\rangle  =  2P^0\sum_{n=0}^\infty C_n ({\mu}^2 z^2){\cal F}_n(-\lambda)\sum_{m=0}^n{\cal B}_{nm}(\mu)\nn\\
	&\hspace{3em}\times\xi^n\int_{-1}^1 dx\ C_m^{3/2}\left({x\over \xi}\right) F(x,\xi,t,\mu)+\ldots\,,
\end{align}
where ${\cal F}_n(-\lambda)$
are partial wave polynomials whose explicit forms are known in the conformal OPE of current-current correlators for the hadronic light-cone DAs~\cite{Braun:2007wv}, ${\cal B}_{nm}$ can be found in~\cite{Braun:2003rp,Mueller:1993hg}, and $\ldots$ denotes the higher-twist contributions $\mathcal{O}\left({M^2/ ({\bar P}^z)^2},{t/ ({\bar P}^z)^2},z^2\Lambda_{\rm QCD}^2\right)$.

Fourier transforming the l.h.s of the above equation to momentum space and invert it order by order in $\alpha_s$, we then obtain the following EFT expansion of the unpolarized quark GPD,
\begin{align}\label{eq:gpdfact}
	&F(x,\xi,t,\mu) \\
	&\hspace{0.5em}=  \int_{-\infty}^\infty {dy\over |\xi|}\bar{C}\left({x\over \xi}, {y\over \xi}, {\mu \over \xi {\bar P}^z }\right){\tilde F}(y,\xi,t, {\bar P}^z,\mu) +\ldots\nn\\
	&\hspace{0.5em}= \int_{-\infty}^\infty {dy\over |y|}C\left({x\over y}, {\xi\over y}, {\mu \over y{\bar P}^z }\right){\tilde F}(y,\xi,t,{\bar P}^z,\mu)+ \ldots\,,\nn
\end{align}
which has been organized following the same spirit as the factorization of PDFs in previous sections. Both forms have been used in the literature~\cite{Ji:2015qla,Xiong:2015nua,Liu:2019urm} with the matching coefficients being related by
\beq
C\left({x\over y}, {\xi\over y}, {\mu \over y{\bar P}^z }\right) = \left|{y\over \xi}\right|\bar{C}\left({x\over \xi}, {y\over \xi}, {\mu \over \xi {\bar P}^z }\right)\,,
\eeq
and $\ldots$ denotes the higher-twist contributions which have the same power-counting as in Eq.~(\ref{eq:ope-twist2}) except that $z^2$ is replaced by $1/(x{\bar P}^z)^2$. For the helicity and transversity quark quasi-GPDs, the factorization formula has the same form as \eq{gpdfact}~\cite{Liu:2019urm}.

The matching coefficient can be obtained by replacing the hadron states in Eqs. (\ref{eq:GPD}) and (\ref{eq:quasiGPD}) with the quark states carrying momentum $p+\Delta/2$ and $p-\Delta/2$ with $p^\mu=(p^0,0,0,p^z)$, and calculating the quark matrix element in perturbation theory.
%
The explicit expression for the $\mathcal O(\alpha_s)$ matching coefficients can be found in~\cite{Liu:2019urm}.
%
An important feature of the result is: The quasi-GPDs do not vanish in all $y$ range, but the collinear singularities only show up in DGLAP and ERBL regions at one-loop. They are exactly the same as those in light-cone GPDs, and thus cancel in the matching coefficient. Moreover, one can derive momentum RGEs for the quasi-GPDs,
which are turned into RGE for the scale dependence of the GPDs by
the matching procedure.


To conclude this subsection, let us make some remarks on the EFT formula for the quark GPD above. First, at zero skewness $\xi=0$, we have
\begin{align} \label{eq:zeroskew}
	F(x,0,t,\mu)&=\!\! \int_{-\infty}^\infty\! {dy\over |y|}C\!\left({x\over y}, 0, {\mu \over y{P}^z }\right)\! {\tilde F}(y,0,t,P^z,\mu)\nn\\
	&\qquad +\ldots\,,
\end{align}
where the matching kernel $C(x/y,0,\mu/(yP^z))$ is exactly the same as the matching coefficient for the quasi-PDF~\cite{Izubuchi:2018srq}, even when $t\neq0$. This can be understood as follows: At zero skewness, both the longitudinal momentum transfer and the energy transfer vanish, the momentum transfer is purely transverse and thus is not affected by Lorentz boost along the longitudinal $z$ direction. As a result, no extra matching related to $t$ is required in the large $P^z$ limit, and the matching remains the same as in the quasi-PDF case. If we take the forward limit $\Delta \to0$, then Eq.~(\ref{eq:zeroskew}) reduces exactly to the EFT expansion formula for the PDF~\cite{Ji:2015qla,Izubuchi:2018srq}.

Second, in the limit $\xi\to1$ and $t\to0$, the quasi-GPD reduces to the quasi-DA that will be discussed in the next subsection, and the corresponding matching kernel also reduces to that for the quasi-DA.


%------------------------------------------------------------------------------
\subsection{Hadronic Distribution Amplitudes}
\label{sec:others-da}
%-----------------------------------------------------------------------------

Within LaMET, the DAs of protons as well as other hadrons can also be extracted from lattice simulations of appropriately chosen quasi-DAs. In this subsection, we show how this can be done in practice. For illustration, we take the leading-twist pion DA as an example. The application to other hadrons~\cite{Chen:2017gck,Wang:2019msf} is analogous.


The leading-twist DA of the pion is the simplest and most extensively studied hadronic DA. It represents the probability amplitude of finding the valence $q\bar q$ Fock state in the pion with the quark carrying a fraction $x$ of the total pion momentum, and is defined as
\begin{align}\label{eq:pionDA}
	\phi_{\pi}(x)&={\frac{1}{if_{\pi}}}\int \frac{d\lambda}{2\pi P^+}e^{-i (x-\frac 1 2)\lambda}\langle 0|O_{\gamma^+\gamma_5}(\lambda n)|\pi(P)\rangle\,,
\end{align}
with normalization $\int_0^1 dx\, \phi_\pi(x)=1$. Here $f_\pi$ denotes the decay constant, and
$O_{\gamma^+\gamma_5}(\lambda n)$ has the same structure as that used in Eq.~(\ref{eq:GPD}) with $\gamma^+$ replaced by $\gamma^+\gamma_5$. The pion DA can be constrained from experimental measurements of, e.g.,
$\gamma \gamma* \rightarrow \pi^0$ from BaBar and Belle~\cite{Aubert:2009mc,Uehara:2012ag},
and then used as an input to test QCD in other measurements such as the pion form factor~\cite{Efremov:1978rn,Farrar:1979aw}.
In the asymptotic limit, it is well known that the pion DA takes the form
$6 x(1-x)$~\cite{Efremov:1978rn,Lepage:1979zb}.
However, how it behaves at lower scales remains under debate (see e.g.~\cite{Chernyak:1981zz}).
%
Calculating the pion DA with controllable systematics in LaMET will be able to shed new lights on its shape and thus on our understanding of pion structure.


Following the same strategy as before, we can access the $x$-dependence of the pion DA by studying the following quasi-DA~\cite{Ji:2015qla,Zhang:2017bzy}
\begin{align}\label{eq:quasiDA}
	\tilde\phi_{\pi}(y, P^z)&={\frac{1}{if_{\pi}}}\int \frac{d\lambda}{2\pi P^z}e^{i (y-\frac 1 2) \lambda}\langle 0|{O}_{\gamma^z\gamma_5}(z)|\pi(P)\rangle\,,
\end{align}
The longitudinally and transversely polarized vector meson quasi-DAs can be defined analogously by replacing $\gamma^z\gamma_5$ in the above equation with $\gamma^0$, $\gamma^z\gamma_\perp$, respectively~\cite{Liu:2018tox}.


The quark bilinear operators defining quasi-DAs follow the same renormalization pattern as those defining the quasi-PDFs or quasi-GPDs.
In the literature, the Wilson-line mass-subtraction scheme was used in the first LaMET calculations of the meson DAs~\cite{Zhang:2017bzy,Chen:2017gck}, and the RI/MOM scheme has been adopted in more recent works~\cite{Zhang:2020gaj}.

The LaMET expression for DAs takes the following form in the $\MS$ scheme~\cite{Ji:2015qla,Liu:2018tox}
\begin{align}\label{eq:matching2}
	\phi_\pi(x,\mu) &=\int_{-\infty}^\infty dy\, C_{\pi}\left(x,y,P^z/\mu\right) \tilde\phi_\pi(y,P^z,\mu)+\ldots\,.
\end{align}
The matching coefficient  for the quasi-DAs can be obtained by replacing the meson state $| \pi(P)\rangle$ in Eqs. (\ref{eq:pionDA}) and (\ref{eq:quasiDA}) with the lowest Fock state $| q(y P) \bar q((1-y)P)\rangle$ and calculating the quark matrix elements, where $y P$ and $(1-y)P$ are the momenta of the quark $q$ and anti-quark $\bar q$, respectively. Its one-loop results have been calculated in both $\MS$ and RI/MOM schemes~\cite{Liu:2018tox}, which agrees with matching coefficient for the quasi-GPDs~\cite{Ji:2015qla,Xiong:2015nua,Liu:2019urm} in \eq{gpdfact} with the replacement of $\xi\to 1/(2y-1)$, $x/\xi\to 2x-1$, and the external momentum $p^z$ to $p^z/2$.

Apart from the LaMET approach in momentum space,
the shape of the pion DA has also been studied using equal-time current-current correlation in coordinate space approach~\cite{Bali:2017gfr,Bali:2018spj},
\begin{align}
	&\langle 0 | T\Big\{J_\mu\Big(\frac{z}{2}\Big) J_\nu\Big(-\frac{z}{2}\Big)\Big\}| \pi^0(P)\rangle \nn\\
	& \hspace{3em}= \frac{2i\, f_\pi}{3 \pi^2 z^4} \epsilon_{\mu\nu\alpha\beta} P^\alpha z^\beta \Phi_\pi(\lambda, z^2)\,,
	\label{eq:pigg-in-z-space}
\end{align}
where $\Phi_\pi(\lambda, z^2)$ can be factorized as
\begin{align} \label{twist_decomposition_heuristic}
	\Phi_\pi(\lambda, z^2) &= C_2 (\lambda,z^2\mu^2, x)\otimes \phi_\pi(x,\mu)+\cdots\,.
\end{align}
Here the matching coefficient $C_2$ depends  on the choice of the currents. Its explicit expression can be found in~\cite{Bali:2018spj}.
The above factorization is controlled by ${\mathcal O}(z^2\Lambda_{\rm QCD}^2)$, with power corrections denoted by ``$\cdots$''. In~\cite{Bali:2018spj}, a combined analysis of several current-current correlations has been performed where twist-four contributions were also included using the model estimate in~\cite{Braun:1989iv,Ball:2006wn}. {The leading-twist pion DA was then extracted from a global fit to the data, and the second moment of the pion DA has been fitted with controlled precision}, both of which favor a considerably broader shape than the asymptotic DA at a scale of $2$ GeV.
A large pion momentum is required to access information at large $\lambda$ so that we can extract wider range of $x$ or higher moments of the pion DA~\cite{Bali:2017gfr}.

%------------------------------------------------------------------------------
\subsection{Higher-Twist Distributions}
\label{sec:others-ht}
%------------------------------------------------------------------------------

Higher-twist distributions are quantities of great interest because they describe the coherent quark-gluon correlations in
the proton. In contrast with the leading-twist distributions, our understanding of the higher-twist ones is rather poor.
On one hand, they often depend on more than one parton momentum fractions; on the other hand, there is no physical intuition about what they may look like, in particular, about how they behave asymptotically at small and large $x$~\cite{Braun:2011aw}. There have been studies on the higher-twist distributions in the context of their connection to the DIS structure function, the transverse single-spin asymmetries in various hadron productions, GPDs related to quark and gluon OAM, parton DAs, etc. LaMET will be able to shed new lights by providing a possibility to access them from the lattice.

Higher-twist contributions also appear in LaMET expansion, where
the suppression is provided by powers of the hadron momentum squared.
In all factorizations presented in previous sections, only the leading-twist terms that capture the logarithmic dependence on hadron momentum are taken into account. The higher-twist contributions have been assumed to be small. If the hadron momentum is not sufficiently large compared and/or one is close to the endpoint region ($x\to0$ and $x\to1$), the higher-twist contributions can become non-negligible, whose structure and impact require understanding.


\subsubsection{Higher-twist collinear-parton observables}

Beyond leading-twist, there exist three simplest twist-three quark distributions
$e(x)$, $g_T(x)$ and $h_L(x)$ related to the unpolarized, transversely and longitudinally polarized
proton~\cite{Jaffe:1991ra},
\begin{align}
	e(x) &= \frac{1}{2M}\int \frac{d\lambda}{2\pi}e^{ix\lambda} \\
	&\quad\times \langle PS|\psi^\dagger_+(0)\gamma_0\psi_-(\lambda n)|PS\rangle + {\rm h.c.}  \,, \nn\\
	g_T(x)& = \frac{1}{2M}\int \frac{d\lambda}{2\pi}e^{ix\lambda}  \\
	&\quad\times \langle PS_\perp|\psi^\dagger_+(0)\gamma_0\gamma_\perp\gamma_5\psi_-(\lambda n)|PS_\perp\rangle + {\rm h.c.} \,,\nn\\
	h_L(x)&  =\frac{1}{2M}\int \frac{d\lambda}{2\pi} e^{ix\lambda} \\
	&\quad\times \langle PS_z|\psi^\dagger_+(0)\gamma_0\gamma_5\psi_-(\lambda n)|PS_z\rangle + {\rm h.c.}  \,,\nn
	\label{eq:twist3}
\end{align}
where we have again employed the decomposition of quark fields $\psi=\psi_+ +\psi_-$ in Sec.~\ref{sec:lc} and the light-cone gauge $A^+=0$, and ``h.c.'' stands for Hermitian conjugate.

The twist-three distributions can contribute as leading effects in certain experimental observables. For example, $g_T(x)$ and $h_L(x)$ can be measured as the leading effects in the longitudinal-transverse spin asymmetry in polarized Drell-Yan process.

Since $\psi_-$ is a non-dynamical component depending on $\psi_+$, all the above distributions can be shown to be related to more complicated quark-gluon correlation functions~\cite{Ji:1990br,Balitsky:1996uh}. A complete set of such correlation functions has been given in~\cite{Qiu:1991pp,Ji:1992eu,Ji:2001bm,Kang:2008ey}, where the quark-gluon correlations
in a transversely-polarized proton take the following form
\begin{align}
	&T_q(x_1, x_2)=\frac{1}{(P^+)^2}\int\frac{d\lambda d\zeta}{(2\pi)^2}e^{i\lambda x_1+i\zeta(x_2-x_1)}\\
	&\hspace{2em}\times\langle PS_\perp|\bar\psi(0)\gamma^+ \epsilon^{+-S_\perp i}g F^{+i}(\zeta n)\psi(\lambda n)|PS_\perp\rangle,\nn\\
	&T_{\Delta q}(x_1, x_2)=\frac{1}{(P^+)^2}\int\frac{d\lambda d\zeta}{(2\pi)^2}e^{i\lambda x_1+i\zeta(x_2-x_1)}\\
	&\hspace{2em}\times\langle PS_\perp|\bar\psi(0)i\gamma^+\gamma_5 S^i_\perp g F^{+i}(\zeta n)\psi(\lambda n)|PS_\perp\rangle\,.\nn
\end{align}
There are also ones in an unpolarized and longitudinally-polarized proton.
Generalizing to off-forward kinematics, the resulting twist-three GPDs are also related to quark
and gluon OAM contribution to the proton spin~\cite{Hatta:2012cs,Ji:2012ba}.

One can also define twist-four distributions in a similar way as in Eq. (\ref{eq:twist3})
by using $\psi_-$ for both quark fields. More general twist-four distributions
will involve three light-cone variables, which will contribute to, e.g.,
$1/Q^2$ term in DIS~\cite{Jaffe:1982pm,Ellis:1982cd,Jaffe:1991ra,Ji:1993ey}.

In principle, all the above higher-twist distributions, as well as others that have not been listed here, can be
computed using the LaMET approach by choosing appropriate quasi-LF correlations. For example, the first exploratory lattice calculation of $g_T(x)$ has been done in~\cite{Bhattacharya:2020cen}, which will be discussed in \sec{lattice-other_distributions}. However, extra complications are expected due to their complex structure. For instance, the lightcone zero modes that do not enter in dealing with leading-twist distributions come into play here. Recently, one of the authors has shown how to study the properties of these zero modes from lattice simulations in LaMET~\cite{Ji:2020baz}. In addition, the higher-twist distributions will have a more complex mixing pattern~\cite{Ji:1990br,Balitsky:1996uh}. Thus,
their matching from the corresponding quasi distributions must take into account such
mixings, making them more challenging than calculating the twist-two PDFs.  One-loop studies of the matching for twist-three disributions have been carried out in~\cite{Bhattacharya:2020xlt,Bhattacharya:2020jfj}.


\subsubsection{Higher-twist contributions to quasi-PDFs}

Let us turn to the power suppressed higher-twist contributions appearing in the extraction of leading-twist quark PDFs using LaMET. Such contributions have two distinct origins. To understand them, let us recall the OPE for the quasi-LF correlation in \eq{ope}. For simplicity, we ignore the renormalization here.
%
Recovering the leading-twist quark PDF requires removing the contributions of both trace terms in that equation.
The trace terms on the r.h.s. of \eq{ope}, which lead to contributions suppressed by powers of $M^2/(P^z)^2$, are known as kinematic power contributions or target mass corrections. In DIS, they can be accounted for by changing the scaling variable $x$ to the Nachtmann variable~\cite{Nachtmann:1973mr}. In the case of LaMET, it behaves slightly differently, as shown in the following. The second type of power corrections come from the trace terms in the operators on the r.h.s. of \eq{ope-tq}, and in general leads to contributions of $\mathcal O(\Lambda^2_{\rm{QCD}}/(P^z)^2)$ . These are genuine higher-twist contributions that involve multi-parton correlations, sometimes also known as dynamical higher-twist contributions. The target mass corrections have been computed to all orders in $M^2/(P^z)^2$ for the quark quasi-PDFs in~\cite{Chen:2016utp,Radyushkin:2017ffo}. The genuine higher-twist contributions have been investigated using two different approaches~\cite{Chen:2016utp,Braun:2018brg}.


According to~\cite{Chen:2016utp}, the $M^2/(P^z)^2$ corrections can be obtained from the ratio
%
\begin{align}
	K_m  &\equiv
	\frac{n_{(\mu_1}\cdots n_{\mu_m)}
		P^{\mu_1}\cdots P^{\mu_m}}{n_{\mu_1}\cdots n_{\mu_m}
		P^{\mu_1}\cdots P^{\mu_m}} = \sum_{i=0}^{i_\text{max}} C_{m-i}^i c^i   \label{Kn}\,,
\end{align}
where $i_\text{max}=(m-\text{Mod}[m,2])/2$,
$C$ is the binomial function and
$c=-n^2 M^2/4\left(n\cdot P\right)^2 = M^2/4 (P^z)^2$
with $n^\mu = (0,0,0,-1)$ and $n \cdot P=P^z$.

Plugged into the tree-level OPE formula in \eq{ope}, the above factors can then be converted to the following relation between unpolarized PDF and quasi-PDF~\cite{Chen:2016utp}
\begin{align}
	q(x)&=\sqrt{1+4c}\sum_{n=0}^\infty \frac{(4c)^n}{f_+^{2n+1}}\Big[(1+(-1)^n)\tilde q\Big(\frac{f_+^{2n+1}x}{2(4c)^n}\Big)\nn\\
	&\qquad+(1-(-1)^n)\tilde q\Big(\frac{-f_+^{2n+1}x}{2(4c)^n}\Big)\Big],
\end{align}
where $f_{+}=\sqrt{1+4c}+ 1$. It is worth noting that quark number conservation is preserved in the above result. The target mass corrections for the longitudinally and transversely polarized quasi-PDFs can be derived analogously.

The trace part on the r.h.s. of \eq{ope-tq} is a genuine higher-twist effect. One may try to construct a non-local form of the higher-twist operators from OPE. The leading trace term, which is a twist-four effect, has been studied in~\cite{Chen:2016utp} (see also~\cite{Balitsky:1987bk}) and shown to give rise to a twist-four PDF
\begin{align}\label{tildeq-twist4}
	&{q}_\text{4}(x,P^z)
	=  \int_{-\infty}^\infty \!\frac{d\lambda}{8\pi P^z}\,
	\Gamma_0\left(-ix\lambda\right)
	\left\langle P\left\vert  O_\text{tr}(z)\right\vert P\right\rangle ,
\end{align}
with
\begin{align}
	&O_\text{tr}(z) =
	\int_0^z \!dz_1\, \bar{\psi}(0) \Big[
	\Gamma^\nu W\left(0,z_1\right) D_\nu W\left(z_1,z\right)\\
	&\hspace{-.5em}+ \int_0^{z_1} \!\! \! dz_2\,  n \cdot  \Gamma\,
	W\left(0,z_2\right) D^\nu W\left(z_2,z_1\right) D_\nu
	W\left(z_1,z\right) \Big] \psi(z n)\,,\nn
\end{align}
where one has $\Gamma^{\mu}=\gamma^{\mu}, \gamma^{\mu}\gamma^5, \gamma^\perp\gamma^{\mu}\gamma^5$ for the unpolarized, helicity and transversity PDFs, respectively. $\Gamma_0$ is the incomplete Gamma function
\begin{equation}
	\Gamma_0\left(-ix\right) = \int_0^1 \frac{dt}{t} e^{ix/t} \,.
\end{equation}
The above twist-four contribution needs to be removed from the quasi-PDF to recover the leading-twist PDF. It also provides a possibility for practical computations on the lattice. However, as a multi-parton correlation involving more gauge links and covariant derivatives, its lattice computation is rather challenging and has not been carried out in any existing work yet.


Another approach that has been used to estimate power corrections related to quark quasi-PDFs is the
renormalon model (see~\cite{Beneke:1998ui} for a comprehensive review).
%
It is based on the observation that the perturbative expansion of the matching coefficient for the quasi-PDF diverges factorially with the loop order, implying that it is only well defined up to a power accuracy. This is known as the renormalon ambiguity, which must be cancelled by terms in the higher-twist contributions.


In~\cite{Braun:2018brg}, it was shown that the cancellation of renormalon ambiguity requires that the leading higher-twist or twist-four contribution takes the following form
\begin{align}
	\label{eq:Q4-cutoff}
	{q}_4(y,P^z,\mu) = \mu^2 \int_{-1}^1 \!\frac{dx}{|x|} D\Big(\frac{y}{x}\Big) q(x,\mu) + q'_4(y,P^z,\mu)\,,
\end{align}
where the first term on the r.h.s. cancels the renormalon ambiguity from the leading-twist matching coefficient, and $q'_4$ depends on $\mu$ at most logarithmically. Since the first term is to merely cancel similar contributions in the matching coefficient, it does not contribute to any physical observable. The renormalon model of power corrections \cite{Beneke:1995pq,Dokshitzer:1995qm,Dasgupta:1996hh,Dasgupta:1996ki,Beneke:1997sr,Braun:2004bu} is based on the assumption that, by replacing $\mu$ with a suitable nonperturbative scale, this contribution reflects the order and the functional form of actual power-suppressed contributions. This was known as ``ultraviolet dominance'' in~\cite{Braun:1995tp,Beneke:1998ui,Beneke:2000kc}. Under this assumption, we obtain the following estimate,
\begin{equation}
	{q}_4(y,P^z,\mu) = \kappa
	\Lambda_{\rm QCD}^2 \int_{-1}^1 \!\frac{dx}{|x|} D\Big(\frac{y}{x}\Big)\, q(x,\mu)\,,
	\label{UV_Lambda_int}
\end{equation}
where $\kappa$ is a dimensionless coefficient of $\mathcal O(1)$ that cannot be fixed within theory and remains a free parameter.


A detailed analysis~\cite{Braun:2018brg} showed that for the quasi-PDF we have
\begin{align}
	\label{t4-Qparallel}
	& {q}_4(y,P^z) =\frac{\kappa  \Lambda_{\rm QCD}^2}{y^2(1-y) (P^z)^2}\\
	&\hspace{-.2em}\times \hspace{-.2em}(1-y)\!\Big[\!\!
	\int_{|y|}^1 \!\!\frac{dx}{x} \Big[\frac{x^2}{(1-x)_+}\!-\!2x^2\Big]q\Big(\frac{y}{x}\Big)\!
	+\! 2q(y)\! -\! |y|q'(y) \Big]\,,\nn
\end{align}
where the first term in the integral was reproduced in a recent analysis of the renormalon effects in the quasi-PDF~\cite{Liu:2020rqi}. As one can see, the second row vanishes as $q(y)$ when $y\to1$ if $\lim_{y\to1}q(y)\sim (1-y)^a$ with $a>0$.
This gives an estimate of the twist-four contribution on the r.h.s. of Eq.~(\ref{eq:fact_0}), which  implies that the higher-twist contributions are enhanced as $1/y^2$ and $1/(1-y)$ for $y\sim 0$ and $y\sim 1$, respectively. Similar analysis can also be done for the pseudo-PDF.
The above result can be used as a way to model the twist-four contribution with $\kappa$ as the only parameter.

\subsection{Orbital Angular Momemntum of Partons in the Proton}
\label{sec:spin}
%------------------------------------------------------------------------------

Over the past three decades, much experimental and theoretical
work has been done on the origin and structure of proton spin, which has been covered in depth in the review articles
~\cite{Filippone:2001ux,Bass:2004xa,Aidala:2012mv,Leader:2013jra,Ji:2016djn,Deur:2018roz,Ji:2020ena}.

In addition to the spin-dependent PDFs and TMDs, the GCPOs---in particular the GPDs---also play an important role in understanding the spin structure of the proton. Since GPDs describe the correlation between the transverse position and longitudinal momentum of quarks and gluons inside the proton, they offer a unique channel to study the orbital angular momentum (OAM) from experiments.

There are two widely known definitions of OAM in literature. One is the kinetic OAM in the gauge-invariant and frame-independent  sum rule for the proton spin~\cite{Ji:1996ek,Ji:1996nm}, which is related to the first moment of twist-two GPDs and can be calculated from the form factors of the symmetric QCD energy-momentum tensor. A review of the lattice calculations of kinetic OAM can be found in~\cite{Ji:2020ena}. The other definition, which has a clear partonic interpretation in comparison to the kinetic OAM, is the canonical OAM in the naive partonic sum rule~\cite{Jaffe:1989jz} based on the free-field form of the QCD angular momentum,
\begin{align}\label{eq:jmam}
	\vec{J} &= \int d^3\xi\ \psi^\dagger \frac{\vec{\Sigma}}{2} \psi + \int d^3\xi\ \psi^\dagger\left[\vec{\xi}\times (-i\vec{\nabla} )\right]\psi \nn\\
	&\ \ \ \ \ + \int d^3\xi\ \vec{E}\times\vec{A} + \int d^3\xi\ E^i\ \left(\vec{\xi}\times\vec{\nabla}\right) A^{i} \ ,
\end{align}
where $i$ is the spatial Lorentz index. Except for the first one,
the other three operators are gauge dependent, and their matrix elements are
generally frame dependent. In high-energy scattering, there is one frame and
gauge that are special: the IMF and light-front gauge, $A^+=0$.
Therefore, the naive partonic sum rule for proton spin can be expressed as~\cite{Jaffe:1989jz}
\beq\label{eq:jm}
\frac{1}{2}= \frac{1}{2}\Delta \Sigma(\mu)  + l^z_q(\mu) + \Delta G(\mu) + l^z_g(\mu) \,,
\eeq
where $l^z_q(\mu)$ and $l^z_g(\mu)$ are the canonical
OAM of the quark and gluon partons, respectively.
Both $l_q^z$ and $l_g^z$ can be related to twist-three GPDs~\cite{Hatta:2011ku,Ji:2012ba,Hatta:2012cs}, which can be accessed through spin-asymmetries in hard exclusive processes~\cite{Ji:2016jgn,Hatta:2016aoc,Bhattacharya:2017bvs,Bhattacharya:2018lgm} (see the recent review~\cite{Ji:2020ena}).

To fully understand the partonic spin structure of the proton,
one also needs to determine the quark and gluon canonical OAM, $l^z_q$ and $l^z_g$.
LaMET allows extraction of $l_q^z$ and $l_g^z$ from lattice calculation in the same way as the gluon helicity that was reviewed in \sec{deltaG}.

The quasi-partonic OAM operators can be chosen as the free-field operators fixed in gauges that
belong to the universality class~\cite{Hatta:2013gta}.
Their matrix elements $\tilde l^z_q$ and $\tilde l^z_g$ can be calculated from the off-forward matrix elements of the relevant energy-momentum tensors~\cite{Zhao:2015kca}, for example,
\begin{align}
	\tilde l^z_q(2S^z) &= \lim_{\Delta\to0} \epsilon^{ij}{\partial \over \partial i\Delta^i}\langle P'S| \psi^\dagger (0) i\partial^j \psi(0)|PS\rangle \Big|_{\vec{\nabla}\cdot \vec{A}=0}\,.
\end{align}
where the kinematics is the same as \eq{kinematic_variable}.

Along with $\Delta G$,
$\tilde l^z_q$ and $\tilde l^z_g$ can be matched to the partonic
quantities defined in the Jaffe-Manohar sum rule through the factorization formulas,
\begin{align}\label{eq:oammatching}
	\tilde l^z_q(P^z,\mu)  &= P_{qq}  l^z_q(\mu) + P_{gq} l^z_g(\mu) \nn\\
	&\qquad+ p_{qq} \Delta \Sigma(\mu) + p_{gq}\Delta G(\mu) + ...\, ,\\
	\tilde l^z_g(P^z,\mu)  &= P_{qg}  l^z_q(\mu) + P_{gg} l^z_g(\mu) \nn\\
	&\qquad+ p_{qg} \Delta \Sigma(\mu) + p_{gg}\Delta G(\mu) + ...\, ,
\end{align}
where $\cdots$ are power corrections suppressed by the momentum $P^z$, and the one-loop matching coefficients in front of each term on the r.h.s. have been calculated in the Coulomb gauge~\cite{Ji:2014lra}. Since the quasi-partonic operators are gauge-variant and need to be fixed in a particular gauge, they can mix with new operators that are not allowed by
Lorentz or gauge symmetries. For example,
the gauge-dependent potential angular momentum $\psi^\dagger(\vec{r}\times \vec{A})\psi$
comes into play~\cite{Wakamatsu:2014zza,Ji:2015sio}. Such mixings must be taken into account in lattice renormalization  to have a controlled calculation of the canonical OAM.

Apart from the above approach, it has also been proposed to calculate the ratio of $l_q^z$ and the valence quark number from the derivatives of off-forward matrix elements of staple-shaped quark Wilson line operators~\cite{Engelhardt:2017miy}, whose definition can be found in \eq{quasi_TMD} below. The first lattice calculations with this approach have been carried out in~\cite{Engelhardt:2017miy,Engelhardt:2019lyy}, which shows different size of effects between the kinetic and canonical OAM. {For systematic improvement in this calculation, one should include the matching of such matrix elements to the physical $l_q^z$ in the limit when the transverse separation of the quark fields approaches zero.}

For the transverse polarization, it is natural to define a twist-two
partonic transverse angular momentum density of quarks~\cite{Hoodbhoy:1998yb,Ji:2012sj,Ji:2020hii,Ji:2020ena},
\begin{equation}
	J^{q}_\perp (x) = x\left[q(x)+ E_q(x)\right]/2,
\end{equation}
and similarly for the gluons, where $q(x)$ is the unpolarized quark/antiquark distributions,
and $E_{q,g}(x)$ is the GPDs defined earlier in this section.
Thus to get a simple partonic picture of the
proton transverse spin from the first principles,
it is important to calculate the GPD $E(x)$ using LaMET.

%------------------------------------------------------------------------------
